% Slide 3 — Dataset
\begin{frame}{Dataset}
  \begin{columns}[T,onlytextwidth]
    \column{0.52\textwidth}
    \begin{block}{Dataset snapshot}
      \small
      \begin{tabular}{@{}ll@{}}
        \toprule
        Classes & \NClasses{} (Glioma, Meningioma, Pituitary) \\
        Total images & \NTotal{} (\NPerClass{} per class) \\
        Split & \NTrain{} train / \NTest{} test (80/20, stratified) \\
        Image size & \ImgSize{} grayscale \\
        \bottomrule
      \end{tabular}
    \end{block}

    \begin{block}{Evaluation choice}
      \small We report macro-averaged metrics (especially macro-F1) so each tumor class
      contributes equally to the final score. With balanced classes, improvements reflect
      better separability rather than class reweighting.
    \end{block}

    \column{0.46\textwidth}
    \centering
    \includegraphics[width=\linewidth]{class_balance_train.png}
    \vspace{0.2em}

    {\scriptsize Class distribution (train split).}
  \end{columns}
\end{frame}

\note{
Balanced dataset with 400 training and 100 test samples per class.
This allows fair comparison across tumor types.
}
